\section{Who Splits the Money Today}
\label{sec:related}

We group earlier work by who decides how the money is split, and when (Table~\ref{tab:related}).

\textbf{An optimizer outside the agents.} BAMAS~\cite{yang2026bamas} chooses the models with an integer program and how
the agents work together with reinforcement learning, before the task starts. MasRouter~\cite{yue2025masrouter} sends
each query through a learned controller. ATOM~\cite{zhao2026atom} predicts how hard a query is and sets a budget for
it, and MALBO~\cite{sabbatella2025malbo} and SC-MAS~\cite{zhao2026scmas} tune who is on the team.
ZEBRA~\cite{hamri2026zebra} splits dollars across the steps of a fixed pipeline with a solver before anything runs, and
its language-model baseline loses to a knapsack solver. Mixture-of-Models~\cite{pecerskis2026mixtureofmodels} runs
preset team profiles. In all of these, the agents doing the work do not hold the money. \bench{} includes this setup as
a baseline (a knapsack solver over measured success rates, \S\ref{sec:purse-allocators}).

\textbf{Routing and cascades.} Model routers pick one model per query. RouteLLM~\cite{ong2025routellm} learns from
human preference data, Hybrid LLM~\cite{ding2024hybrid} picks between a small and a large model by predicted
difficulty, and FrugalGPT~\cite{chen2024frugalgpt} tries models from cheap to expensive until an answer is accepted.
RouterBench~\cite{hu2024routerbench} scores routers on recorded outputs against the best possible router, the same
method \bench{}'s replay uses. A router prices one query at a time. \bench{} prices a project whose jobs share one
budget, so money spent on one job is money another job cannot use.

\textbf{Agents that fund agents.} Agent runtimes released in 2026 let a model fund the agents it starts. Automaton sends
credits to a child agent running the same runtime~\cite{conway2026automaton}, and wayland-core's \texttt{delegate}
tool gives a sub-agent a dollar limit out of what its parent has left, charged to the sub-agent and the agents below
it~\cite{ferrox2026waylandcore}. AgentLance~\cite{liu2026markets} lets a manager model set enforced budgets for
subtasks, one level deep, with a market choosing the workers, and evaluates the whole system on benchmarks. \sys{}
v1~\cite{pai2025hashiru} came before these, with a CEO that hires and fires agents under cost and memory budgets, one
level deep. None of them measures how well the agents divide the money.

\textbf{Budget awareness in one agent.} BATS~\cite{liu2026budgetaware} tells a tool-using agent how much budget it has
left. BAGEN~\cite{lin2026bagen} finds that strong agents are not always budget-aware, and are too optimistic about tasks
they cannot finish with the money left. $R^3$-Bench~\cite{wang2026r3bench} has one model split its own output tokens or
actions across six problems of hidden difficulty, and scores it against the best possible split, computed offline from
how each problem responds to more effort. \bench{} computes its best split the same way. EcoAgent-Bench~\cite{wu2026ecoagent}
prices an agent's actions (lookups, searches, a stronger model, a human), AllocBench~\cite{wang2026allocbench} measures
investment in reusable tools, and CostBench~\cite{liu2026costbench} measures planning for the lowest cost and
replanning when tools fail or prices change. In all of them, one agent spends on itself.

\textbf{Local and cloud models, and agent economies.} Minions~\cite{narayan2025costefficient} and
PAPILLON~\cite{li2025papillon} pair a local model with a cloud model under a fixed protocol. Virtual agent
economies~\cite{tomasev2025virtual} and Magentic Marketplace~\cite{bansal2025magentic} study markets between agents.
Payment protocols (AP2~\cite{google2026agent}, the A2A x402 extension~\cite{googleagenticcommerce2025a2a}) move the
money but do not decide how to split it. \sys{} v2 uses the open protocols these build on, MCP for
tools~\cite{modelcontextprotocol2026specification} and A2A for agents~\cite{a2aproject2026agent2agent}.

\begin{table}[!htb]
\centering
\caption{Who splits the money, and when. Like the 2026 agent runtimes, \sys{} v2 lets agents hold budgets. Unlike
them, it caps each call before it is made and is compared with the best possible split.}
\label{tab:related}
\footnotesize
\begin{tabular}{@{}p{4.2cm}p{3.6cm}p{2.6cm}p{3.6cm}@{}}
\toprule
\tablehead Line of work & Who splits the money & When & Compared with the best split \\
\midrule
BAMAS, ZEBRA, MasRouter & optimizer & before the task & no \\
RouteLLM, FrugalGPT & router, cascade & per query & best router \\
Automaton, wayland-core & the model & during & no \\
AgentLance & manager model & per subtask & no \\
$R^3$-Bench, EcoAgent & one agent, on itself & during & yes \\
\textbf{\sys{} v2} & \textbf{the agents, to depth 2} & \textbf{during, each call capped} & \textbf{yes} \\
\bottomrule
\end{tabular}
\end{table}

\sys{} v2 shares agent-held budgets with these runtimes, and borrows $R^3$-Bench's way of computing the best split, but
it differs in two ways. (1) \emph{Limits before the call.} A hire sets money (and memory) aside that firing refunds, and
each call's output is capped at what its agent has left before the call is made. wayland-core's sub-agent, by
contrast, stops itself only after it has gone over its limit. (2) \emph{Compared with the best split.} To our
knowledge, it is the first system whose agent-held budgets, across a hiring tree, are compared with a solver, an
up-front plan and fixed teams at the same budget.
